\documentclass[]{article}

\usepackage{amsmath}
\usepackage{multirow}
\usepackage{natbib}
\usepackage{graphicx} 


\begin{document}

Metastable materials, which reside in local minima on the potential energy landscape above the thermodynamic ground state, represent a rich frontier for discovering novel functionalities. These non-equilibrium phases can host unique electronic, catalytic, and optical properties that are inaccessible in their stable counterparts. However, a fundamental challenge often hinders their practical realization: the energetic penalty of metastability frequently correlates with mechanical fragility. This inherent instability poses a significant barrier to the synthesis, isolation, and ultimate application of these promising materials, making their discovery a resource-intensive and often serendipitous process.Transition-metal dichalcogenides (TMDs) exemplify this paradigm. While their stable semiconducting polymorphs are well-established platforms for next-generation electronics, their metastable phases, such as the metallic 1T and 1T' structures, are highly coveted for applications ranging from electrocatalysis to quantum information. The synthesis of these functional metastable TMDs requires a delicate balance, as many theoretically promising candidates are too mechanically soft to be experimentally isolated. The vast combinatorial space defined by different transition metals and chalcogens makes an exhaustive search for mechanically robust candidates via either first-principles calculations or experimental trial-and-error computationally intractable.To accelerate the discovery of resilient metastable materials, we introduce a machine learning framework designed to efficiently navigate the TMD chemical space and identify promising candidates for synthesis. Our approach bypasses the direct and computationally demanding calculation of elastic properties for every hypothetical structure. Instead, we train a model to predict a key descriptor of mechanical behavior—Pugh's ratio, $G/K$, the ratio of the shear modulus $G$ to the bulk modulus $K$—from a small set of fundamental electronic and structural features. By learning the complex, non-linear relationships between these accessible descriptors and mechanical robustness, our model can rapidly screen large databases of materials to pinpoint those with a high likelihood of being both functional and durable.In this work, we develop a predictive model founded on physically interpretable descriptors, such as the density of states at the Fermi level, which provides insight into the electronic origins of mechanical softening. We rigorously evaluate our model's ability to generalize to new chemical families using a stringent leave-one-metal-group-out cross-validation scheme, a critical test for any data-driven materials discovery effort aiming to extrapolate beyond its training data. By coupling our model's predictions with robust uncertainty quantification, we screen a library of 112 theoretical metastable TMDs. The result is a high-confidence viability map that balances predicted mechanical robustness against thermodynamic accessibility, providing a targeted and prioritized roadmap for the experimental synthesis of novel and resilient functional materials.

\end{document}
                